\newpage
\begin{problem}[Warmup]{(a)}
    In each of the following cases, show that the given element is algebraic over the given base field, by finding an appropriate polynomial with the element as a root.
    \begin{enumerate}
        \item[(a)] $\sqrt[3]{2}$ is algebraic over $\mathbb{Q}$.
        \item[(b)] $\sqrt[3]{2}$ is algebraic over $\mathbb{R}$.
        \item[(c)] $1 + \sqrt{2}$ is algebraic over $\mathbb{Q}$.
        \item[(d)] $\sqrt{1 + \sqrt{2}}$ is algebraic over $\mathbb{Q}$.
        \item[(e)] $1 + 2i$ is algebraic over $\mathbb{Q}$.
        \item[(f)] $\pi + i$ is algebraic over $\mathbb{R}$.
    \end{enumerate}
\end{problem}
\unsolved

\begin{problem}[Warmup]{(b)}
    If $F$ is a field, is $F$ an extension field of $F$? If so, does $F$ have any transcendental elements over $F$?
\end{problem}
\unsolved

\begin{problem}[Warmup]{(c)}
    What extension field do we get if we apply Kronecker's Theorem 42.1 to a degree 1 polynomial?
\end{problem}
\unsolved

\begin{problem}[Warmup]{(d)}
    Recall the Fundamental Theorem of Algebra 9.1, which asserts that all non-constant polynomials in $\mathbb{C}[x]$ can be factored completely into linear factors. What does this say about finding elements outside $\mathbb{C}$ that are algebraic over $\mathbb{C}$?
\end{problem}
\unsolved

\begin{problem}[Warmup]{(e)}
    Give examples of the following, or else explain why no such example exists:
    \begin{enumerate}
        \item[(a)] An element in $\mathbb{C}$ transcendental over $\mathbb{R}$.
        \item[(b)] An element in $\mathbb{C}$ transcendental over $\mathbb{Q}$.
        \item[(c)] An element in $\mathbb{R}$, but not in $\mathbb{Q}(\sqrt{2})$, which is algebraic over $\mathbb{Q}(\sqrt{2})$.
    \end{enumerate}
\end{problem}
\unsolved

\begin{problem}[Warmup]{(f)}
    For each of the following fields, identify its characteristic, and determine its prime subfield: \[\mathbb{C}, \quad \mathbb{Q}, \quad \mathbb{Q}(\sqrt{2}), \quad \mathbb{Z}_{11}, \quad \mathbb{Q}[x]/\langle x^3 - 2 \rangle, \quad \mathbb{Z}_3[x]/\langle x^2 + x + 2 \rangle.\]
\end{problem}
\unsolved

\begin{problem}[Warmup]{(g)}
    ``Every polynomial of degree at least 1 has a root.'' Discuss the truth of this statement.
\end{problem}
\unsolved

\begin{problem}[Warmup]{(h)}
    Why is $\cos 20^\circ$ algebraic over $\mathbb{Q}$? \emph{Hint:} Find the appropriate result in Chapter 39.
\end{problem}
\unsolved

\begin{problem}[Warmup]{(i)}
    Give some examples of real numbers that are algebraic over $\mathbb{Q}$, but not constructible.
\end{problem}
\unsolved

\begin{problem}{1}
    Let $F$ be a field and $f \in F[x]$ a polynomial of degree at least 1.
    Prove that there exists a field extension $E$ of $F$ in which $f$ can be
    factored into linear factors. Such a field extension is called a
    \textbf{splitting extension for $f$}.
\end{problem}
\solved

\begin{proof}
    Let $n = \deg f \geq 1$.
    By Kronecker's
    Theorem there exists an extension field $E_1$ of $F$ such that
    $\alpha_1 \in E_1$ and $f(\alpha_1) = 0$. That is, by the Factor Theorem,
    in the field $E_1$, $f$ can be factored
    \[
        f = (x - \alpha_1)\, g(x), \qquad g \in E_1[x],
    \]
    where $\deg(g) = \deg(f) - 1$. Now repeat this for $g(x)$, applying
    Kronecker's Theorem to $g$ over the field $E_1$ (as long as
    $\deg g \geq 1$), and so on. Since the degree drops by one at each step,
    after finitely many steps we obtain a finite list of extension fields
    \[
        F \subseteq E_1 \subseteq E_2 \subseteq \cdots \subseteq E_n
    \]
    where, in $E_n[x]$,
    \[
        f = c\,(x - \alpha_1)(x - \alpha_2)\cdots(x - \alpha_n),
    \]
    with $\alpha_i \in E_n$ and $c \in F$ the leading coefficient of $f$.
    This $E_n$ is an extension field of $F$. 
    We have thus counstructed an extension field of $F$ wherein $f$ can be 
        factored into linear factors as required.
\end{proof}

\begin{problem}{2}
    In the proof of Kronecker's Theorem 42.1, we prove that the function
    $\psi : F \to F[x]/\langle p \rangle$ is a one-to-one ring homomorphism.
    Show that $\psi$ is onto if and only if the irreducible polynomial $p$
    is linear.
\end{problem}
\solved

\begin{proof}
    Suppose $\psi$ is onto. Then $\langle p \rangle + x = \psi(c) = \langle p \rangle + c$
    for some $c \in F$, so $x - c \in \langle p \rangle$, so $p$ divides $x - c$.
    Since $x - c \neq 0$, this gives $\deg p \leq \deg(x - c) = 1$, and since $p$ is
    irreducible, $\deg p \geq 1$. Thus $\deg p = 1$, so $p$ is linear.

    Conversely, suppose the irreducible polynomial $p$ is linear, say
    $p = ax + b$ with $a, b \in F$ and $a \neq 0$.
    Take $f \in F[x]/\langle p \rangle$, so $f = \langle p \rangle + g$ for some $g \in F[x]$.
    By the division algorithm, $g = qp + r$ where $q \in F[x]$ and either $r = 0$
    or $\deg r < \deg p = 1$, so $r \in F$.
    Since $qp \in \langle p \rangle$,
    \[f = \langle p \rangle + (qp + r) = \langle p \rangle + r = \psi(r).\]
    Thus $\psi$ is onto.
\end{proof}

\begin{problem}{3}
    If $R$ is a commutative ring with unity, and $R$ is a subring of a field
    $F$, we can speak of the elements in $F$ that are \textbf{algebraic} over
    $R$, meaning simply those elements of the field which are roots of
    polynomials from $R[x]$.
    \begin{enumerate}
        \item[(a)] Explain why all rational numbers are algebraic over $\mathbb{Z}$.
        \item[(b)] Prove that the set of elements of $\mathbb{R}$ that are
        algebraic over $\mathbb{Z}$ is the same as the set of elements of
        $\mathbb{R}$ which are algebraic over $\mathbb{Q}$.
    \end{enumerate}
\end{problem}
\solved

\begin{proof}
    Take $c = a/b \in \mathbb{Q}$ where $a, b \in \mathbb{Z}$, $b \ne 0$ and $f = bx - a \in \mathbb{Z}[x]$.
    Then $f(c) = f(a/b) = b(a/b) - a = a - a = 0$ thus $c$ is algebraic over $\mathbb{Z}$.

    Let $r \in \mathbb{R}$ and $c_i = a_i/b_i \in \mathbb{Q}$ where $a_i, b_i \in \mathbb{Z}$, $b_i \neq 0$,
    and not all $c_i = 0$.
    Let $B = b_0 b_1 \cdots b_n$, then
    \[\sum c_i r^i = 0 \iff \sum \frac{a_i}{b_i} r^i = 0 \iff \sum a_i \frac{B}{b_i} r^i = B \cdot 0 \iff \sum a_i \frac{B}{b_i} r^i = 0,\]
    where each $a_i B / b_i \in \mathbb{Z}$ and not all are $0$.
    Thus the set of elements of $\mathbb{R}$ that are
        algebraic over $\mathbb{Z}$ is the same as the set of elements of
        $\mathbb{R}$ which are algebraic over $\mathbb{Q}$.
\end{proof}

\begin{problem}{4}
    Suppose that $\alpha \in \mathbb{C}$ is an algebraic number over
    $\mathbb{Q}$, and $r \in \mathbb{Q}$. Prove that $\alpha + r$ and
    $r\alpha$ are also algebraic over $\mathbb{Q}$. (Actually, this exercise
    is a special case of Exercise 44.10, where you will show that the set
    $\mathbb{A}$ of complex numbers algebraic over $\mathbb{Q}$ is a field,
    called the \textbf{field of algebraic numbers}. Why is our exercise a
    special case?)
\end{problem}
\solved
 
\begin{proof}
    Since $\alpha \in \mathbb{C}$ is an algebraic number over
    $\mathbb{Q}$ there exists a nonzero $f \in \mathbb{Q}[x]$ such that $f(\alpha) = 0$.
    We can suppose $r \ne 0$ (since if $r = 0$ then $r \alpha = 0$, which is a root of $x \in \mathbb{Q}[x]$).
    Let $g(x) = f(r^{-1}x) \in \mathbb{Q}[x]$, which is nonzero since it has the same degree as $f$
    and its leading coefficient is $r^{-n}$ times that of $f$, where $n = \deg f$. Then
    \[g(r \alpha) = f(r^{-1} r \alpha) = f(\alpha) = 0.\]
    Thus $r \alpha$ is algebraic over $\mathbb{Q}$.

    Similarly, let $h(x) = f(x - r) \in \mathbb{Q}[x]$, which is nonzero since it has the same
    degree and leading coefficient as $f$. Then
    \[h(\alpha + r) = f(\alpha + r - r) = f(\alpha) = 0.\]
    Thus $\alpha + r$ is algebraic over $\mathbb{Q}$.

    This is a special case of Exercise 44.10 since every $r \in \mathbb{Q}$ is algebraic over
    $\mathbb{Q}$ (it is a root of $x - r$), so $r, \alpha \in \mathbb{A}$, and since $\mathbb{A}$ is a
    field it is closed under addition and multiplication, so $\alpha + r, r\alpha \in \mathbb{A}$.
\end{proof}

\begin{problem}{5}
    Prove that $\sin 1^\circ$ is algebraic over $\mathbb{Q}$.
\end{problem}
\unsolved

\begin{problem}{6}
    Let $F$ be a field with characteristic $p$. Prove that the function
    $\iota : \mathbb{Z}_p \to F$ discussed in Section 42.3 is a one-to-one
    ring homomorphism.
\end{problem}
\unsolved

\begin{proof}
    Let $x, y \in \mathbb{Z}_p$.
    Then 
    \[\iota([x + y]) = \]
    and 
\end{proof}

\begin{problem}{7}
    \begin{enumerate}
        \item[(a)] For a field $F$ with characteristic zero, we defined in
        Section 42.3 a one-to-one ring homomorphism
        $\iota : \mathbb{Q} \to F$. Prove that $\iota(\mathbb{Q})$ is the
        smallest subfield of $F$.
        \item[(b)] For a field $F$ with characteristic $p$, we defined in
        Section 42.3 a one-to-one ring homomorphism
        $\iota : \mathbb{Z}_p \to F$. (You completed the proof of this in
        Exercise 6.) Prove that $\iota(\mathbb{Z}_p)$ is the smallest
        subfield of $F$.
    \end{enumerate}
\end{problem}
\unsolved

\begin{problem}{8}
    Let $K$ be a field, and consider the set of all subfields of $K$ (this
    set may well have infinitely many elements). Consider the intersection
    of all these subfields. Prove that the result is a field, by checking
    explicitly the field axioms.
\end{problem}
\unsolved
